\documentclass{article}
\usepackage{amsmath,amssymb}
\usepackage[vocab={polynomials}]{leanify}
\newtheorem{lemma}{Lemma}
\begin{document}

\section{Notation}

We write $\operatorname{lc}(p)$ for the leading coefficient of $p$.
We write $a_k(p)$ for the $k$th coefficient of $p$.
We write $N(p)$ for the number of real roots of $p$.
We write $Z(p)$ for the set of real roots of $p$.

\section{Polynomials}

\begin{lemma}\label{product}
Let $p$ and $q$ be nonzero polynomials. Then $\deg(pq)=\deg p+\deg q$ and $\operatorname{lc}(pq)=\operatorname{lc}(p)\operatorname{lc}(q)$.
\end{lemma}

\begin{lemma}\label{roots}
Let $n$ be a positive integer. For every polynomial $p$ of degree $n$, $N(p)\le n$.
\end{lemma}

\begin{lemma}\label{odd}
Let $p$ be a cubic polynomial. Then there is $x\in\mathbb{R}$ such that $p(x)=0$.
\end{lemma}

\begin{lemma}\label{monic}
Let $p$ be a monic polynomial with integer coefficients, and let $r$ be a rational number with $p(r)=0$. Then $r$ is an integer.
\end{lemma}

\begin{lemma}\label{irreducible}
Let $p$ be a quadratic polynomial with $N(p)=0$. Then $p$ is irreducible over $\mathbb{R}$.
\end{lemma}

\begin{lemma}\label{coefficients}
Let $p$ be a polynomial. Then $a_0(p)=p(0)$.
\end{lemma}

\begin{lemma}\label{set}
Let $p$ be a nonzero polynomial and let $a\in Z(p)$. Then $p(a)=0$.
\end{lemma}

\end{document}
